% Options for packages loaded elsewhere
\PassOptionsToPackage{unicode}{hyperref}
\PassOptionsToPackage{hyphens}{url}
\PassOptionsToPackage{dvipsnames,svgnames,x11names}{xcolor}
%
\documentclass[
  10pt,
  a4paper,
]{article}
\usepackage{amsmath,amssymb}
\usepackage{setspace}
\usepackage{iftex}
\ifPDFTeX
  \usepackage[T1]{fontenc}
  \usepackage[utf8]{inputenc}
  \usepackage{textcomp} % provide euro and other symbols
\else % if luatex or xetex
  \usepackage{unicode-math} % this also loads fontspec
  \defaultfontfeatures{Scale=MatchLowercase}
  \defaultfontfeatures[\rmfamily]{Ligatures=TeX,Scale=1}
\fi
\usepackage{lmodern}
\ifPDFTeX\else
  % xetex/luatex font selection
  \setmainfont[]{Liberation Serif}
  \setsansfont[]{Liberation Sans}
  \setmonofont[]{DejaVu Sans Mono}
\fi
% Use upquote if available, for straight quotes in verbatim environments
\IfFileExists{upquote.sty}{\usepackage{upquote}}{}
\IfFileExists{microtype.sty}{% use microtype if available
  \usepackage[]{microtype}
  \UseMicrotypeSet[protrusion]{basicmath} % disable protrusion for tt fonts
}{}
\makeatletter
\@ifundefined{KOMAClassName}{% if non-KOMA class
  \IfFileExists{parskip.sty}{%
    \usepackage{parskip}
  }{% else
    \setlength{\parindent}{0pt}
    \setlength{\parskip}{6pt plus 2pt minus 1pt}}
}{% if KOMA class
  \KOMAoptions{parskip=half}}
\makeatother
\usepackage{xcolor}
\usepackage[margin=22mm]{geometry}
\setlength{\emergencystretch}{3em} % prevent overfull lines
\providecommand{\tightlist}{%
  \setlength{\itemsep}{0pt}\setlength{\parskip}{0pt}}
\setcounter{secnumdepth}{-\maxdimen} % remove section numbering
\ifLuaTeX
\usepackage[bidi=basic]{babel}
\else
\usepackage[bidi=default]{babel}
\fi
\babelprovide[main,import]{korean}
\ifPDFTeX
\else
\babelfont{rm}[]{Liberation Serif}
\fi
% get rid of language-specific shorthands (see #6817):
\let\LanguageShortHands\languageshorthands
\def\languageshorthands#1{}
\usepackage{xeCJK}
\setCJKmainfont{Noto Serif CJK KR}
\setCJKsansfont{Noto Sans CJK KR}
\setCJKmonofont{Noto Sans Mono CJK KR}
\xeCJKsetup{CJKspace=true}
\usepackage{xurl}
\usepackage{fvextra}
\usepackage{fancyhdr}
\AtBeginDocument{\DefineVerbatimEnvironment{verbatim}{Verbatim}{breaklines=true,breakanywhere=true,fontsize=\small}}
\setlength{\emergencystretch}{4em}
\setlength{\headheight}{15pt}
\pagestyle{fancy}\fancyhf{}
\fancyhead[L]{\small 중성자별 구조와 비방사 진동}
\fancyfoot[C]{\thepage}
\renewcommand{\headrulewidth}{0.3pt}
\setcounter{tocdepth}{2}
\ifLuaTeX
  \usepackage{selnolig}  % disable illegal ligatures
\fi
\usepackage{bookmark}
\IfFileExists{xurl.sty}{\usepackage{xurl}}{} % add URL line breaks if available
\urlstyle{same}
\hypersetup{
  pdftitle={A06 · TOV: 중심에서 유효 표면까지},
  pdflang={ko-KR},
  colorlinks=true,
  linkcolor={Maroon},
  filecolor={Maroon},
  citecolor={Blue},
  urlcolor={Blue},
  pdfcreator={LaTeX via pandoc}}

\title{A06 · TOV: 중심에서 유효 표면까지}
\usepackage{etoolbox}
\makeatletter
\providecommand{\subtitle}[1]{% add subtitle to \maketitle
  \apptocmd{\@title}{\par {\large #1 \par}}{}{}
}
\makeatother
\subtitle{과제 큐시트}
\author{}
\date{}

\begin{document}
\maketitle

\setstretch{1.12}
\textbf{선행:} A05　 \textbf{독립 작업:} 6--10시간 (계획치)

\subsection{TOV를 이미 풀었는데 왜 중심과 표면을 다시
읽는가?}\label{tovuxb97c-uxc774uxbbf8-uxd480uxc5c8uxb294uxb370-uxc65c-uxc911uxc2ecuxacfc-uxd45cuxba74uxc744-uxb2e4uxc2dc-uxc77duxb294uxac00}

중심 시작점을 바꾸는 실험과 표면 cutoff를 바꾸는 실험을 나란히 놓자.
각각 어떤 근사 또는 문제의 정의를 바꾸는지 먼저 구별한다.

먼저 한두 문장으로 답한 뒤 아래 설명을 읽자. 첫 예측이 틀렸거나 아직
모르겠다고 적는 것은 감점하지 않는다.

\subsection{필요한 관계와 읽을
위치}\label{uxd544uxc694uxd55c-uxad00uxacc4uxc640-uxc77duxc744-uxc704uxce58}

\(m^{\prime}=4\pi r^2\epsilon/c^2\),
\(P^{\prime}=-(\epsilon+P)\nu^{\prime}\),
\(\nu^{\prime}=G(m+4\pi r^3P/c^2)/[c^2r^2(1-2Gm/(rc^2))]\).

\begin{itemize}
\tightlist
\item
  \path|notes/chapters/ch03_tov.md|
\item
  \path|notes/equations/tov.md|
\item
  \path|learner/src/structure/ns_tov.f90|
\end{itemize}

\subsection{계산으로
옮기기}\label{uxacc4uxc0b0uxc73cuxb85c-uxc62euxae30uxae30}

기본 경로는 제공 함수의 호출과 해석이다. 아래 \texttt{learner} 시험은
같은 물리 함수를 직접 작성한 경우에 실행한다. 제공 예제를 검사할 때에는
명시적으로 \texttt{calculation}을 선택한다.

\begin{enumerate}
\def\labelenumi{\arabic{enumi}.}
\tightlist
\item
  r=0에 RHS를 직접 넣지 않고 \(m\sim r^3\), \(P-P_c\sim r^2\)의 계수를
  구한다.
\item
  q=ln(Pstart/P)에서 dr/dq의 부호를 유도한다.
\item
  양의 Pcut에서 정지하고 lapse의 적분상수를 같은 기준으로 맞춘다.
\item
  r0, ODE 허용오차, Pcut 변화는 별개의 실험으로 기록한다.
\end{enumerate}

\newpage

\subsection{실행해서 무엇을 볼
것인가}\label{uxc2e4uxd589uxd574uxc11c-uxbb34uxc5c7uxc744-uxbcfc-uxac83uxc778uxac00}

패키지 루트에서 실행한다. \texttt{-\/-plan}은 경로만 확인하며 계산
성공을 뜻하지 않는다.

\begin{verbatim}
python tools/run_assignment.py A06 --plan
python tools/run_assignment.py A06 --run \
  --workspace learner --output ../A06_result.json
\end{verbatim}

\subsection{결과를
설명하기}\label{uxacb0uxacfcuxb97c-uxc124uxba85uxd558uxae30}

중심 급수, 적분 오차, 표면 선택은 서로 다르다. 질량·반경의 일치만 보고
lapse 기준까지 같다고 판단하지 않는다.

\textbf{남길 증거.} 중심급수 유도; M,R 및 μexp(ν) profile; 세 오차원에
대한 별도 비교표

\textbf{판단 기준.} 2점 예: 중심 시작 오차, 적분 오차, 표면 모델
민감도를 한 숫자로 합치지 않는다. 최초 예측의 정오보다, 결과를 본 뒤
이유를 고치고 다른 입력에도 적용하는지를 본다.

\subsection{막혔을 때}\label{uxb9c9uxd614uxc744-uxb54c}

\textbf{힌트 1.} m은 kg이며 Gm/c²와 혼동하지 않는다.

\textbf{힌트 2.} Pcut을 바꾸는 실험에서는 같은 IVP의 step refinement가
아니라 정의한 절단 문제가 달라진다.

\subsection{다음 만남에 다시
묻기}\label{uxb2e4uxc74c-uxb9ccuxb0a8uxc5d0-uxb2e4uxc2dc-uxbb3buxae30}

다음 회차에는 자신의 TOV 한 행에서 진동에 필요한 계수를 찾는다. 없는
양은 0으로 채우지 말고 어디서 계산할지 설명한다. 먼저 짧게 답한 뒤
본문과 자신의 기록을 열어 확인한다.

\textbf{범위.} 양압 cutoff는 정확한 진공 표면과의 접합이 아니다.
회전·자기장·탄성 지각은 범위 밖이다.

\end{document}
